\section{Curves}
\label{section-curves}

\noindent
In the Stacks project we will use the following as our definition
of a curve.

\begin{definition}
\label{definition-curve}
Let $k$ be a field. A {\it curve} is a variety of dimension $1$ over $k$.
\end{definition}

\noindent
Two standard examples of curves over $k$ are the affine line $\mathbf{A}^1_k$
and the projective line $\mathbf{P}^1_k$. The scheme $X = \Spec(k[x, y]/(f))$
is a curve if and only if $f \in k[x, y]$ is irreducible.

\medskip\noindent
Our definition of a curve has the same problems as our definition of a
variety, see the discussion following Definition \ref{definition-variety}.
Moreover, it means that every curve comes with a specified field of definition.
For example $X = \Spec(\mathbf{C}[x])$ is a curve over $\mathbf{C}$
but we can also view it as a curve over $\mathbf{R}$. The scheme
$\Spec(\mathbf{Z})$ isn't a curve, even though the schemes $\Spec(\mathbf{Z})$
and $\mathbf{A}^1_{\mathbf{F}_p}$ behave similarly in many respects.

\begin{lemma}
\label{lemma-proper-minus-point}
Let $X$ be a separated, irreducible scheme of dimension $> 0$ over a field $k$.
Let $x \in X$ be a closed point. The open subscheme $X \setminus \{x\}$
is not proper over $k$.
\end{lemma}

\begin{proof}
Since $X$ is irreducible, $U = X \setminus \{x\}$ is not closed in $X$.
In particular, the immersion $U \to X$ is not proper.
By Morphisms, Lemma \ref{morphisms-lemma-image-proper-scheme-closed}
(here we use $X$ is separated), $U \to \Spec(k)$ is not proper either.
\end{proof}

\begin{lemma}
\label{lemma-dim-1-quasi-projective}
Let $X$ be a separated finite type scheme over a field $k$.
If $\dim(X) \leq 1$ then $X$ is H-quasi-projective over $k$.
\end{lemma}

\begin{proof}
By Proposition \ref{proposition-dim-1-noetherian-separated-has-ample}
the scheme $X$ has an ample invertible sheaf $\mathcal{L}$.
By Morphisms, Lemma \ref{morphisms-lemma-quasi-projective-finite-type-over-S}
we see that $X$ is isomorphic to a locally
closed subscheme of $\mathbf{P}^n_k$ over $\Spec(k)$. This is
the definition of being H-quasi-projective over $k$, see
Morphisms, Definition \ref{morphisms-definition-quasi-projective}.
\end{proof}

\begin{lemma}
\label{lemma-dim-1-proper-projective}
Let $X$ be a proper scheme over a field $k$.
If $\dim(X) \leq 1$ then $X$ is H-projective over $k$.
\end{lemma}

\begin{proof}
By Lemma \ref{lemma-dim-1-quasi-projective} we see that $X$ is a
locally closed subscheme of $\mathbf{P}^n_k$ for some field $k$.
Since $X$ is proper over $k$ it follows that $X$ is a closed subscheme
of $\mathbf{P}^n_k$
(Morphisms, Lemma \ref{morphisms-lemma-image-proper-scheme-closed}).
\end{proof}

\begin{lemma}
\label{lemma-dim-1-projective-completion}
Let $X$ be a separated scheme of finite type over $k$.
If $\dim(X) \leq 1$, then there exists an open immersion
$j : X \to \overline{X}$ with the following properties
\begin{enumerate}
\item $\overline{X}$ is H-projective over $k$, i.e., $\overline{X}$
is a closed subscheme of $\mathbf{P}^d_k$ for some $d$,
\item $j(X) \subset \overline{X}$ is dense and scheme
theoretically dense,
\item $\overline{X} \setminus X = \{x_1, \ldots, x_n\}$
for some closed points $x_i \in \overline{X}$.
\end{enumerate}
\end{lemma}

\begin{proof}
By Lemma \ref{lemma-dim-1-quasi-projective} we may assume $X$ is a
locally closed subscheme of $\mathbf{P}^d_k$ for some $d$. Let
$\overline{X} \subset \mathbf{P}^d_k$ be the scheme theoretic image
of $X \to \mathbf{P}^d_k$, see Morphisms, Definition
\ref{morphisms-definition-scheme-theoretic-image}.
The description in
Morphisms, Lemma \ref{morphisms-lemma-quasi-compact-immersion}
gives properties (1) and (2).
Then $\dim(X) = 1 \Rightarrow \dim(\overline{X}) = 1$ for example by
looking at generic points, see
Lemma \ref{lemma-dimension-locally-algebraic}.
As $\overline{X}$ is Noetherian, it then
follows that $\overline{X} \setminus X = \{x_1, \ldots, x_n\}$
is a finite set of closed points.
\end{proof}

\begin{lemma}
\label{lemma-reduced-dim-1-projective-completion}
Let $X$ be a separated scheme of finite type over $k$.
If $X$ is reduced and $\dim(X) \leq 1$, then there exists
an open immersion $j : X \to \overline{X}$ such that
\begin{enumerate}
\item $\overline{X}$ is H-projective over $k$, i.e., $\overline{X}$
is a closed subscheme of $\mathbf{P}^d_k$ for some $d$,
\item $j(X) \subset \overline{X}$ is dense and scheme
theoretically dense,
\item $\overline{X} \setminus X = \{x_1, \ldots, x_n\}$
for some closed points $x_i \in \overline{X}$,
\item the local rings $\mathcal{O}_{\overline{X}, x_i}$
are discrete valuation rings for $i = 1, \ldots, n$.
\end{enumerate}
\end{lemma}

\begin{proof}
Let $j : X \to \overline{X}$ be as in
Lemma \ref{lemma-dim-1-projective-completion}.
Consider the normalization $X'$ of $\overline{X}$ in
$X$. By Lemma \ref{lemma-relative-normalization-finite}
the morphism $X' \to \overline{X}$ is finite.
By Morphisms, Lemma \ref{morphisms-lemma-finite-projective}
$X' \to \overline{X}$ is projective. By Morphisms, Lemma
\ref{morphisms-lemma-projective-over-quasi-projective-is-H-projective}
we see that $X' \to \overline{X}$ is H-projective.
By Morphisms, Lemma \ref{morphisms-lemma-H-projective-composition}
we see that $X' \to \Spec(k)$ is H-projective.
Let $\{x'_1, \ldots, x'_m\} \subset X'$ be the inverse
image of $\{x_1, \ldots, x_n\} = \overline{X} \setminus X$.
Then $\dim(\mathcal{O}_{X', x'_i}) = 1$ for all $1 \leq i \leq m$.
Hence the local rings $\mathcal{O}_{X', x'}$
are discrete valuation rings by
Morphisms, Lemma
\ref{morphisms-lemma-relative-normalization-normal-codim-1}.
Then $X \to X'$ and $\{x'_1, \ldots, x'_m\}$ is as desired.
\end{proof}

\begin{lemma}
\label{lemma-dim-1-projective-completion-after-insep}
Let $X$ be a separated scheme of finite type over $k$ with $\dim(X) \leq 1$.
Then there exists a commutative diagram
$$
\xymatrix{
\overline{Y}_1 \amalg \ldots \amalg \overline{Y}_n \ar[rd] &
Y_1 \amalg \ldots \amalg Y_n \ar[r]_-\nu \ar[d] \ar[l]^j &
X_{k'} \ar[r] \ar[d] &
X \ar[d]^f \\
& \Spec(k'_1) \amalg \ldots \amalg \Spec(k'_n) \ar[r] &
\Spec(k') \ar[r] &
\Spec(k)
}
$$
of schemes with the following properties:
\begin{enumerate}
\item $k'/k$ is a finite purely inseparable extension of fields,
\item $\nu$ is the normalization of $X_{k'}$,
\item $j$ is an open immersion with dense image,
\item $k'_i/k'$ is a finite separable extension for $i = 1, \ldots, n$,
\item $\overline{Y}_i$ is smooth, projective, geometrically irreducible
dimension $\leq 1$ over $k'_i$.
\end{enumerate}
\end{lemma}

\begin{proof}
As we may replace $X$ by its reduction, we may and do assume $X$ is reduced.
Choose $X \to \overline{X}$ as in
Lemma \ref{lemma-reduced-dim-1-projective-completion}.
If we can show the lemma for $\overline{X}$, then the lemma follows
for $X$ (details omitted).
Thus we may and do assume $X$ is projective.

\medskip\noindent
Choose $k'/k$ finite purely inseparable such that the normalization of
$X_{k'}$ is geometrically normal over $k'$, see
Lemma \ref{lemma-finite-extension-geometrically-normal}.
Denote $Y = (X_{k'})^\nu$ the normalization; for properties
of the normalization, see Section \ref{section-normalization}.
Then $Y$ is geometrically regular as normal and regular are the same
in dimension $\leq 1$, see
Properties, Lemma \ref{properties-lemma-normal-dimension-1-regular}.
Hence $Y$ is smooth over $k'$ by
Lemma \ref{lemma-geometrically-regular-smooth}.
Let $Y = Y_1 \amalg \ldots \amalg Y_n$ be the decomposition
of $Y$ into irreducible components.
Set $k'_i = \Gamma(Y_i, \mathcal{O}_{Y_i})$.
These are finite separable extensions of $k'$ by
Lemma \ref{lemma-proper-geometrically-reduced-global-sections}.
The proof is finished by Lemma \ref{lemma-baby-stein}.
\end{proof}

\begin{lemma}
\label{lemma-regular-point-on-curve}
Let $k$ be a field. Let $X$ be a curve over $k$. Let $x \in X$ be a closed
point. We think of $x$ as a (reduced) closed subscheme of $X$ with sheaf
of ideals $\mathcal{I}$. The following are equivalent
\begin{enumerate}
\item $\mathcal{O}_{X, x}$ is regular,
\item $\mathcal{O}_{X, x}$ is normal,
\item $\mathcal{O}_{X, x}$ is a discrete valuation ring,
\item $\mathcal{I}$ is an invertible $\mathcal{O}_X$-module,
\item $x$ is an effective Cartier divisor on $X$.
\end{enumerate}
If $k$ is perfect or if $\kappa(x)$ is separable over $k$,
these are also equivalent to
\begin{enumerate}
\item[(6)] $X \to \Spec(k)$ is smooth at $x$.
\end{enumerate}
\end{lemma}

\begin{proof}
Since $X$ is a curve, the local ring $\mathcal{O}_{X, x}$ is a Noetherian
local domain of dimension $1$ (Lemma \ref{lemma-dimension-locally-algebraic}).
Parts (4) and (5) are equivalent by definition and are equivalent to
$\mathcal{I}_x = \mathfrak m_x \subset \mathcal{O}_{X, x}$ having one generator
(Divisors, Lemma \ref{divisors-lemma-effective-Cartier-in-points}).
The equivalence of (1), (2), (3), (4), and (5) therefore follows from
Algebra, Lemma \ref{algebra-lemma-characterize-dvr}. The final statement
follows from Lemma \ref{lemma-dense-smooth-open-variety-over-perfect-field}
in case $k$ is perfect. If $\kappa(x)/k$ is separable, then the equivalence
follows from Algebra, Lemma \ref{algebra-lemma-separable-smooth}.
\end{proof}

\begin{remark}
\label{remark-smooth-projective-compactification}
Let $k$ be a field. Let $X$ be a regular curve over $k$. By
Lemmas \ref{lemma-regular-point-on-curve} and
\ref{lemma-reduced-dim-1-projective-completion}
there exists a nonsingular projective curve $\overline{X}$
which is a compactification of $X$, i.e., there exists an
open immersion $j : X \to \overline{X}$ such that the complement
consists of a finite number of closed points. If $k$ is perfect,
then $X$ and $\overline{X}$ are smooth over $k$ and
$\overline{X}$ is a smooth projective compactification of $X$.
\end{remark}

\noindent
Observe that if an affine scheme $X$ over $k$ is proper over $k$
then $X$ is finite over $k$ (Morphisms, Lemma
\ref{morphisms-lemma-finite-proper}) and hence has
dimension $0$
(Algebra, Lemma \ref{algebra-lemma-finite-dimensional-algebra} and
Proposition \ref{algebra-proposition-dimension-zero-ring}).
Hence a scheme of dimension $> 0$ over $k$ cannot be both affine and
proper over $k$. Thus the possibilities in the following lemma are
mutually exclusive.

\begin{lemma}
\label{lemma-curve-affine-projective}
Let $X$ be a curve over $k$. Then either $X$ is an affine scheme or $X$
is H-projective over $k$.
\end{lemma}

\begin{proof}
Choose $X \to \overline{X}$ with
$\overline{X} \setminus X = \{x_1, \ldots, x_r\}$ as in
Lemma \ref{lemma-reduced-dim-1-projective-completion}.
Then $\overline{X}$ is a curve as well.
If $r = 0$, then $X = \overline{X}$ is H-projective over $k$.
Thus we may assume $r \geq 1$ and our goal is to show that
$X$ is affine. By Lemma \ref{lemma-open-in-affine-curve-affine}
it suffices to show that $\overline{X} \setminus \{x_1\}$ is affine.
This reduces us to the claim stated in the next paragraph.

\medskip\noindent
Let $X$ be an H-projective curve over $k$. Let $x \in X$ be a closed point
such that $\mathcal{O}_{X, x}$ is a discrete valuation ring. Claim:
$U = X \setminus \{x\}$ is affine. By Lemma \ref{lemma-regular-point-on-curve}
the point $x$ defines an effective Cartier divisor of $X$.
For $n \geq 1$ denote $nx = x + \ldots + x$ the $n$-fold sum, see
Divisors, Definition \ref{divisors-definition-sum-effective-Cartier-divisors}.
Denote $\mathcal{O}_{nx}$ the structure sheaf of $nx$ viewed as a
coherent module on $X$.
Since every invertible module on the local scheme $nx$ is trivial
the first short exact sequence of
Divisors, Remark \ref{divisors-remark-ses-regular-section}
reads
$$
0 \to \mathcal{O}_X
\xrightarrow{1} \mathcal{O}_X(nx) \to \mathcal{O}_{nx} \to 0
$$
in our case. Note that $\dim_k H^0(X, \mathcal{O}_{nx}) \geq n$.
Namely, by Lemma \ref{lemma-chi-tensor-finite} we have
$H^0(X, \mathcal{O}_{nx}) = \mathcal{O}_{X, x}/(\pi^n)$ where $\pi$ in
$\mathcal{O}_{X, x}$ is a uniformizer and the powers $\pi^i$ map to
$k$-linearly independent elements in $\mathcal{O}_{X, x}/(\pi^n)$
for $i = 0, 1, \ldots, n - 1$. We have $\dim_k H^1(X, \mathcal{O}_X) < \infty$
by Cohomology of Schemes, Lemma
\ref{coherent-lemma-proper-over-affine-cohomology-finite}.
If $n > \dim_k H^1(X, \mathcal{O}_X)$
we conclude from the long exact cohomology sequence
that there exists an $s \in \Gamma(X, \mathcal{O}_X(nx))$
which is not a section of $\mathcal{O}_X$.
If we take $n$ minimal with this property, then $s$ will
map to a generator of the stalk
$\left(\mathcal{O}_X(nx)\right)_x$ since otherwise it would
define a section of $\mathcal{O}_X((n - 1)x) \subset \mathcal{O}_X(nx)$.
For this $n$ we conclude that $s_0 = 1$ and $s_1 = s$
generate the invertible module $\mathcal{L} = \mathcal{O}_X(nx)$.

\medskip\noindent
Consider the corresponding morphism
$f = \varphi_{\mathcal{L}, (s_0, s_1)} : X \to \mathbf{P}^1_k$
of Constructions, Section \ref{constructions-section-projective-space}.
Observe that the inverse image of $D_{+}(T_0)$ is $U = X \setminus \{x\}$
as the section $s_0$ of $\mathcal{L}$ only vanishes at $x$.
In particular, $f$ is non-constant, i.e., $\Im(f)$ has more
than one point. Hence $f$ must map the generic
point $\eta$ of $X$ to the generic point of $\mathbf{P}^1_k$.
Hence if $y \in \mathbf{P}^1_k$ is a closed point, then $f^{-1}(\{y\})$
is a closed set of $X$ not containing $\eta$, hence finite.
Finally, $f$ is proper\footnote{Namely, a H-projective variety
is a proper variety by Morphisms, Lemma
\ref{morphisms-lemma-projective-is-quasi-projective-proper}.
A morphism of varieties whose source is a proper variety is a proper morphism
by Morphisms, Lemma \ref{morphisms-lemma-image-proper-scheme-closed}.}.
By Cohomology of Schemes, Lemma
\ref{coherent-lemma-proper-finite-fibre-finite-in-neighbourhood}\footnote{One
can avoid using this lemma which relies on the theorem of formal
functions. Namely, $X$ is projective hence it suffices to show
a proper morphism $f : X \to Y$ with finite fibres between quasi-projective
schemes over $k$ is finite. To do this, one chooses an affine open of $X$
containing the fibre of $f$ over a point $y$ using that any finite set of
points of a quasi-projective scheme over $k$ is contained in an affine.
Shrinking $Y$ to a small affine neighbourhood of $y$ one reduces to the
case of a proper morphism between affines. Such a morphism is finite by
Morphisms, Lemma \ref{morphisms-lemma-integral-universally-closed}.}
we conclude that $f$ is finite. Hence $U = f^{-1}(D_{+}(T_0))$
is affine.
\end{proof}

\noindent
The following lemma combined with Lemma \ref{lemma-proper-minus-point}
tells us that given a separated scheme $X$ of dimension $1$ and
of finite type over $k$, then $X \setminus Z$ is affine, whenever the
closed subset $Z$ meets every irreducible component of $X$.

\begin{lemma}
\label{lemma-dim-1-nonproper-affine}
Let $X$ be a separated scheme of finite type over $k$.
If $\dim(X) \leq 1$ and no irreducible component of $X$
is proper of dimension $1$, then $X$ is affine.
\end{lemma}

\begin{proof}
Let $X = \bigcup X_i$ be the decomposition of $X$ into irreducible components.
We think of $X_i$ as an integral scheme (using the reduced induced scheme
structure, see Schemes, Definition
\ref{schemes-definition-reduced-induced-scheme}).
In particular $X_i$ is a singleton (hence affine) or a curve
hence affine by Lemma \ref{lemma-curve-affine-projective}.
Then $\coprod X_i \to X$ is finite surjective and $\coprod X_i$ is affine.
Thus we see that $X$ is affine by
Cohomology of Schemes, Lemma
\ref{coherent-lemma-image-affine-finite-morphism-affine-Noetherian}.
\end{proof}






